\begin{abstract}
\noindent
\newline
Las crisis cambiarias, como el colapso de la libra esterlina en 1992 o la crisis del peso mexicano en 1994, generan costes económicos devastadores y a menudo sorprenden a los mercados a pesar de la existencia de modelos de alerta temprana. La literatura económica tradicional explica las causas fundamentales \textit{ex post}, pero carece de mecanismos precisos para detectar el \textit{timing} exacto del colapso en tiempo real.

Este trabajo tiene como objetivo llenar este vacío aplicando la ``Teoría del Peligro'' (\textit{Danger Theory}) de la inmunología a los mercados financieros. Proponemos que los mercados emiten ``señales de peligro'' internas (intervención del banco central, estrés en reservas) análogas a las señales biológicas, las cuales pueden ser detectadas antes de que aparezcan los síntomas externos (colapso de precios).

Para ello, utilizamos una base de datos híbrida que integra series diarias de tipos de cambio con datos de intervención de alta frecuencia del Banco de Inglaterra y estadísticas macroeconómicas del FMI y la ONS para el Reino Unido (1990-1992) y México (1993-1995).

La metodología empleada consiste en un Algoritmo Genético que entrena una población de agentes para detectar patrones no lineales en estas señales de estrés. A diferencia de los modelos econométricos, estos agentes evolucionan sus propias reglas de decisión sin supuestos funcionales a priori.

Nuestra principal aportación es demostrar que un agente artificial puede predecir el ``Miércoles Negro'' con un día de antelación basándose casi exclusivamente en la intervención y el desempleo, ignorando el precio. Además, validamos la robustez del modelo mediante pruebas sintéticas (``Trampa de Valor'') y transferencia de aprendizaje al caso mexicano, donde una recalibración automática del umbral de pánico permitió evitar falsos positivos políticos y predecir el colapso final.

\vspace{0.5cm}
\noindent \textbf{Keywords:} Crisis cambiarias, Teoría del Peligro, Algoritmos Genéticos, Intervención Cambiaria, Sistemas Inmunes Artificiales.

\noindent \textbf{JEL Classification:} F31, C63, G01.
\end{abstract}
